% =============================================================================
% TCC II — alvo SBC (formato de artigo)
% UFOPA / IEG — Bacharelado em Ciência da Computação
%
% Alvo PRINCIPAL do texto, decidido em 24/08/2026. O irmão main.tex produz o
% mesmo conteúdo no layout do guia do SIBI, exigido pela biblioteca no depósito.
%
% OS CAPÍTULOS INCLUÍDOS ABAIXO SÃO OS MESMOS DO main.tex, na mesma ordem, e é
% assim que tem de continuar: tests/test_conformidade_sbc.py compara a lista de
% \input dos dois alvos e falha se divergirem. Layout duplicado é a decisão;
% prosa duplicada é o defeito que ela não autoriza.
%
% O que este alvo NÃO tem, por não existir no formato de artigo: ficha
% catalográfica, folha de aprovação, sumário, dedicatória, agradecimentos e
% epígrafe. Todos continuam no main.tex, que é o documento que vai à biblioteca.
% =============================================================================

\documentclass[12pt]{article}

\usepackage{sbc-tcc}

% --- Dados do trabalho -------------------------------------------------------
% Ver a nota em texto/tcc1/main-sbc.tex: o formato SBC não leva titulação no
% nome do autor, e a orientadora entra como coautora conforme o Art. 7º, § 2º do
% regulamento do curso.
\title{Sistemas inteligentes aplicados à gestão pública: automação da fase
  preparatória e pesquisa de preços em licitações municipais}

\author{Rafael Nobre de Souza\inst{1}, Socorro Vânia Lourenço Alves\inst{1}}

\address{Instituto de Engenharia e Geociências\\
  Universidade Federal do Oeste do Pará (UFOPA)\\
  Santarém --- PA --- Brasil
  \email{\pendente{endereços institucionais, conferidos antes da entrega}}}

\begin{document}

\maketitle

\begin{resumo}
  \pendente{Resumo definitivo: objetivo, método, resultados obtidos e conclusão.
  No máximo 10 linhas, como o formato SBC exige. Não reaproveitar o resumo do
  alvo do SIBI sem cortar — ele é dimensionado para até 500 palavras.}
\end{resumo}

\begin{abstract}
  \pendente{Tradução do resumo definitivo, feita depois que ele estabilizar.}
\end{abstract}

% --- Elementos textuais -------------------------------------------------------
% Mesma lista, mesma ordem que main.tex. Ver o aviso no cabeçalho.
\section{Introdução}

\pendente{Evoluir a introdução do TCC I, incorporando o que a execução do
projeto revelou. Ao migrar texto, reconferir cada citação: afirmação que era
"esperada" em TCC I precisa virar afirmação sustentada por resultado em TCC II,
ou ser removida.}

\section{Fundamentação teórica}

\pendente{Consolidar a partir dos fichamentos do vault. Toda afirmação não
trivial precisa de fonte rastreável até um documento lido — não até uma
lembrança de leitura.}

\section{Trabalhos relacionados}

\pendente{Resultado do levantamento sistemático, com o quadro comparativo e a
lacuna que este trabalho ocupa.}

\section{Materiais e métodos}

\pendente{Descrever o que foi efetivamente feito, no passado, e não o que se
pretendia fazer. Divergências em relação ao planejado no TCC I são material
legítimo do texto: registre a mudança e o motivo.}

\subsection{Arquitetura da solução}

\pendente{Diagrama e descrição dos componentes.}

\subsection{Modelagem de dados}

\pendente{Modelo conceitual e lógico, incluindo o mecanismo de trilha de
auditoria que permite reproduzir uma pesquisa de preços na data em que foi
realizada.}

\subsection{Integração com dados públicos de contratações}

\pendente{Endpoints consumidos, parâmetros de consulta, tratamento de falhas e
estratégia de cache.}

\subsection{Tratamento estatístico dos preços}

\pendente{Critério de exclusão de valores discrepantes e cálculo dos
parâmetros, com a justificativa normativa de cada escolha.}

\subsection{Geração assistida de minutas}

\pendente{Como os documentos são produzidos e como sua conformidade é
verificada.}

\subsection{Procedimento de avaliação}

\pendente{Participantes, tarefas, métricas e instrumentos. Detalhar o suficiente
para que outra pessoa possa repetir o experimento.}

\section{Resultados e discussão}

\pendente{Apresentar os dados obtidos antes de interpretá-los. Números sem
medição não entram — se um resultado não foi medido, ele não existe nesta
seção.}

\section{Considerações finais}

\pendente{Retomar o problema e os objetivos, dizendo objetivamente quais foram
atingidos e quais não foram.}

\subsection{Contribuições}

\pendente{O que este trabalho entrega que antes não existia.}

\subsection{Limitações}

\pendente{Escrever com honestidade. Banca reconhece limitação omitida, e a
omissão custa mais caro que a limitação.}

\subsection{Trabalhos futuros}

\pendente{Desdobramentos concretos, não lista de desejos.}


% --- Referências --------------------------------------------------------------
\ufopabibliografiasbc{refs}

\end{document}
